% Luc Destombe --- licence CC BY-NC-SA 4.0
% https://creativecommons.org/licenses/by-nc-sa/4.0/deed.fr
% Fichier autonome : compiler avec pdflatex, deux fois (signets, nombre de pages).
\def\lyceeversion{1.2}
\documentclass[a4paper,10pt]{article}
\makeatletter
\newif\iflycee@niveaux
\lycee@niveauxtrue
\newif\iflycee@cm
\newif\iflycee@palatino

\RequirePackage[utf8]{inputenc}
\RequirePackage[T1]{fontenc}
\RequirePackage[french]{babel}
\RequirePackage{etoolbox}

\newcommand{\ShorthandsOff}{\@ifundefined{shorthandoff}{}{\shorthandoff{;:!?}}}
\newcommand{\ShorthandsOn}{\@ifundefined{shorthandon}{}{\shorthandon{;:!?}}}
\ShorthandsOff
\AtBeginDocument{\ShorthandsOn}
\AtBeginEnvironment{tikzpicture}{\ShorthandsOff}

\RequirePackage{geometry}
\RequirePackage{fancyhdr}
\RequirePackage{lastpage}
\RequirePackage{multicol}
\RequirePackage{array}
\RequirePackage{enumitem}
\RequirePackage{xcolor}
\RequirePackage{needspace}

\RequirePackage{amsmath,amssymb}
\iflycee@cm\else\RequirePackage{fourier}\fi
\DeclareMathAlphabet{\mathcal}{OMS}{cmsy}{m}{n}
\renewcommand{\ttdefault}{lmtt}
\begingroup\catcode`\_=\active \catcode`\^=\active
\global\let_\sb \global\let^\sp
\endgroup
\AtBeginDocument{\catcode`\_=12 \mathcode`\_="8000
  \catcode`\^=12 \mathcode`\^="8000 }
\newlength{\retraitcalculs}
\setlength{\retraitcalculs}{2em}

\RequirePackage{tikz}
\usetikzlibrary{arrows.meta,calc,babel,shapes.misc}
\RequirePackage[most]{tcolorbox}
\tcbuselibrary{skins,breakable}

\definecolor{coulDef}{HTML}{1F4E79}
\definecolor{coulProp}{HTML}{7030A0}
\definecolor{coulRem}{HTML}{C55A11}
\definecolor{coulEx}{HTML}{2E7D32}
\definecolor{coulExo}{HTML}{B03A2E}
\definecolor{coulPart}{HTML}{1F4E79}
\definecolor{coulMeth}{HTML}{1B6E4F}
\definecolor{coulCourbe}{HTML}{0B5ED7}
\definecolor{coulCourbeB}{HTML}{C00000}
\definecolor{coulCourbeC}{HTML}{2E7D32}
\definecolor{coulGrille}{HTML}{8C8C8C}
\definecolor{coulRep}{HTML}{1B6E4F}
\definecolor{coulBar}{HTML}{2E7D32}

\newcommand{\R}{\mathbb{R}}
\newcommand{\Rs}{\mathbb{R}^{*}}
\newcommand{\N}{\mathbb{N}}
\newcommand{\Ns}{\mathbb{N}^{*}}
\newcommand{\Z}{\mathbb{Z}}
\newcommand{\Q}{\mathbb{Q}}
\renewcommand{\le}{\leqslant}
\renewcommand{\ge}{\geqslant}
\renewcommand{\approx}{\simeq}
\newcommand{\vect}[1]{\overrightarrow{#1}}
\newcommand{\norme}[1]{\left\lVert #1 \right\rVert}
\newcommand{\intervalle}[2]{\left[#1\,;#2\right]}
\RequirePackage{cancel}
\renewcommand{\CancelColor}{\color{coulExo}}
\newcommand{\fois}[1]{\times{\color{coulExo}#1}}
\newcommand{\barre}[1]{\cancel{#1}}
\newlength{\lycee@eqh}\newlength{\lycee@eqd}
\newsavebox{\lycee@eqa}\newsavebox{\lycee@eqb}\newsavebox{\lycee@eqc}
\newcommand{\lycee@eqcol}[1]{\hbox{$\renewcommand{\arraystretch}{1.35}%
  \begin{array}[t]{@{}>{\displaystyle}r@{}>{\displaystyle{}}l@{}}#1\end{array}$}}
\newcommand{\lycee@eqmesure}[1]{%
  \ifdim\ht#1>\lycee@eqh\setlength{\lycee@eqh}{\ht#1}\fi
  \ifdim\dp#1>\lycee@eqd\setlength{\lycee@eqd}{\dp#1}\fi}
\newcommand{\lycee@eqtrait}{\hspace{1.2em}ou\hspace{1.2em}}
\newcommand{\lycee@eqiff}{$\iff$\hspace{0.8em}}
\newcommand{\eqproduit}[2]{\par\smallskip\noindent C'est une équation produit :
  $A\times B=0 \iff A=0$ ou $B=0$.\par\nopagebreak\smallskip
  \begingroup\setlength{\lycee@eqh}{0pt}\setlength{\lycee@eqd}{0pt}%
  \setbox\lycee@eqa=\lycee@eqcol{#1}\setbox\lycee@eqb=\lycee@eqcol{#2}%
  \lycee@eqmesure\lycee@eqa\lycee@eqmesure\lycee@eqb%
  \noindent\hspace*{\retraitcalculs}\lycee@eqiff\usebox\lycee@eqa\lycee@eqtrait\usebox\lycee@eqb\endgroup\par\smallskip}
\newcommand{\eqproduittrois}[3]{\par\smallskip\noindent C'est une équation produit :
  $A\times B\times C=0 \iff A=0$ ou $B=0$ ou $C=0$.\par\nopagebreak\smallskip
  \begingroup\setlength{\lycee@eqh}{0pt}\setlength{\lycee@eqd}{0pt}%
  \setbox\lycee@eqa=\lycee@eqcol{#1}\setbox\lycee@eqb=\lycee@eqcol{#2}\setbox\lycee@eqc=\lycee@eqcol{#3}%
  \lycee@eqmesure\lycee@eqa\lycee@eqmesure\lycee@eqb\lycee@eqmesure\lycee@eqc%
  \noindent\hspace*{\retraitcalculs}\lycee@eqiff\usebox\lycee@eqa\lycee@eqtrait\usebox\lycee@eqb\lycee@eqtrait\usebox\lycee@eqc\endgroup
  \par\smallskip}

\tikzset{grille/.style={coulGrille,line width=0.4pt}}
\newcommand{\quadrillage}[4]{\draw[grille,step=1] (#1,#2) grid (#3,#4);}
\tikzset{
  croix/.style={cross out,draw,line width=0.6pt,inner sep=0pt,minimum size=4pt},
  croixdroite/.style={croix,rotate=45,line width=0.8pt},
}
\newcommand{\pt}[3]{\path (#1) node[croix]{} node[#2,font=\small,inner sep=2.5pt]{$#3$};}

\newcommand{\lycee@repere}[4]{%
  \draw[grille,step=1] (#1,#2) grid (#3,#4);
  \draw[-{Stealth[length=2mm]},thick] (#1,0)--({#3+0.4},0) node[below left=-1pt]{$x$};
  \draw[-{Stealth[length=2mm]},thick] (0,#2)--(0,{#4+0.4}) node[below left=-1pt]{$y$};
  \node[below left,font=\scriptsize,inner sep=1.5pt] at (0,0) {$0$};}
\newcommand{\lycee@gradx}[1]{\foreach \k/\l in {#1}{%
  \draw (\k,0.07)--(\k,-0.07) node[below,font=\scriptsize,inner sep=1.5pt]{$\l$};}}
\newcommand{\lycee@grady}[1]{\foreach \k/\l in {#1}{%
  \draw (0.07,\k)--(-0.07,\k) node[left,font=\scriptsize,inner sep=1.5pt]{$\l$};}}
\newcommand{\lycee@intff}[2]{\left[#1\,;#2\right]}
\newcommand{\lycee@intfo}[2]{\left[#1\,;#2\right[}
\newcommand{\lycee@intof}[2]{\left]#1\,;#2\right]}
\newcommand{\lycee@intoo}[2]{\left]#1\,;#2\right[}
\newcommand{\lycee@ens}[1]{\left\{#1\right\}}
\AtBeginDocument{%
  \providecommand{\repere}{\lycee@repere}%
  \providecommand{\gradx}{\lycee@gradx}%
  \providecommand{\grady}{\lycee@grady}%
  \providecommand{\Cf}{\mathcal{C}\sb{f}}%
  \providecommand{\Cg}{\mathcal{C}\sb{g}}%
  \providecommand{\intff}{\lycee@intff}%
  \providecommand{\intfo}{\lycee@intfo}%
  \providecommand{\intof}{\lycee@intof}%
  \providecommand{\intoo}{\lycee@intoo}%
  \providecommand{\ens}{\lycee@ens}}

\newlist{questions}{enumerate}{3}
\setlist[questions,1]{label=\arabic*\textdegree),leftmargin=*,itemsep=2pt,topsep=3pt}
\setlist[questions,2]{label=\alph*),leftmargin=*,itemsep=1pt,topsep=2pt}
\setlist[questions,3]{label=\roman*),leftmargin=*,itemsep=1pt,topsep=1pt}
\setlist[itemize]{label=\textbullet,leftmargin=*,itemsep=2pt,topsep=3pt}

\newcommand{\besoinplace}[1]{\par\penalty\@M\begingroup
  \dimen@=#1\relax\dimen@ii=\pagegoal\advance\dimen@ii-\pagetotal
  \ifdim\dimen@>\dimen@ii\ifdim\dimen@ii>\z@\vfil\fi\break\fi\endgroup}

\newcommand{\lycee@pied}{\thepage/\pageref{LastPage}}
\newcommand{\entete}[2]{%
  \pagestyle{fancy}\fancyhf{}%
  \fancyhead[L]{\footnotesize\color{coulPart}\textbf{#1}}%
  \fancyhead[R]{\footnotesize\color{coulPart}\textbf{#2}}%
  \fancyfoot[C]{\footnotesize\lycee@pied}%
  \renewcommand{\headrulewidth}{0.4pt}%
  \renewcommand{\headrule}{\hbox to\headwidth{%
    \color{coulPart!40}\leaders\hrule height \headrulewidth\hfill}}}

\RequirePackage{ccicons}
\newcommand{\lycee@licence}{{\scriptsize\color{gray}%
  \copyright\ Luc Destombe --- \ccbyncsa\ CC BY-NC-SA 4.0}}
\newif\iflycee@licence \lycee@licencetrue
\newcommand{\sanslicence}{\lycee@licencefalse}
\AtBeginDocument{\iflycee@licence\AddToHookNext{shipout/foreground}{%
  \put(0,0){\rlap{\kern\dimexpr1in+\hoffset+\oddsidemargin+\textwidth\relax
    \llap{\raisebox{-\dimexpr1in+\voffset+\topmargin+\headheight+\headsep
      +\textheight+\footskip\relax}{\lycee@licence}}}}}\fi}

\newcommand{\lycee@site}{faire-des-maths.fr}
\newcommand{\lycee@version}{{\scriptsize\color{gray}\lycee@site
  \ifdefined\lyceeversion\ --- Version~\lyceeversion\fi}}
\AtBeginDocument{\iflycee@licence\AddToHookNext{shipout/foreground}{%
  \put(0,0){\rlap{\kern\dimexpr1in+\hoffset+\oddsidemargin\relax
    \raisebox{-\dimexpr1in+\voffset+\topmargin+\headheight+\headsep
      +\textheight+\footskip\relax}{\lycee@version}}}}\fi}

\newcommand{\formulesserrees}{%
  \setlength{\abovedisplayskip}{3pt}\setlength{\belowdisplayskip}{3pt}%
  \setlength{\abovedisplayshortskip}{2pt}\setlength{\belowdisplayshortskip}{3pt}}

\setlength{\parindent}{0pt}

\RequirePackage[implicit=false,hidelinks,pdfencoding=auto,psdextra,bookmarksopen,
  bookmarksopenlevel=1,pdfcreator={},pdfproducer={}]{hyperref}
\RequirePackage{bookmark}
\hypersetup{pdfauthor={Luc Destombe},pdfkeywords={CC BY-NC-SA 4.0}}
\newcounter{lycee@signet}
\newcommand{\signet}[3][]{\stepcounter{lycee@signet}%
  \edef\lycee@dest{\if\relax\detokenize{#1}\relax
    signet.\arabic{lycee@signet}\else#1\fi}%
  \hypertarget{\lycee@dest}{}\bookmark[level=#2,dest=\lycee@dest]{#3}}
\pdfstringdefDisableCommands{%
  \def\R{R}\def\N{N}\def\Z{Z}\def\Q{Q}\def\vect#1{#1}\def\Cf{Cf}\def\Cg{Cg}\def\fois#1{×#1}%
  \def\barre#1{#1}\def\intervalle#1#2{[#1;#2]}\def\intff#1#2{[#1;#2]}%
  \def\textsuperscript#1{#1}\def\dfrac#1#2{#1/#2}\def\frac#1#2{#1/#2}%
  \def\vec#1{#1⃗}}%
\RequirePackage[official]{eurosym}

\geometry{top=0.8cm,bottom=1.3cm,left=1cm,right=1cm,footskip=0.6cm}

\pagestyle{fancy}
\fancyhf{}
\renewcommand{\headrulewidth}{0pt}
\fancyfoot[C]{\footnotesize Page \thepage{} / \pageref{LastPage}}

\newcommand{\titrefiche}[2]{%
  \begin{center}
    \Large\textbf{#1}\\[2pt]
    \large\textbf{#2}\\[2pt]
  \end{center}
  \vspace{0.10cm}}

\setlength{\columnseprule}{0.4pt}
\setlength{\columnsep}{15pt}
\renewcommand{\columnseprulecolor}{\color{black!45}}
\setlength{\multicolsep}{2pt}

\newcounter{partiefiche}
\renewcommand{\thepartiefiche}{\Roman{partiefiche}}
\newif\ifapresbandeau
\newcommand{\lycee@bandeau}[1]{%
  \par\penalty-600
  \vspace{0.08cm}%
  \noindent\colorbox{coulPart}{%
    \parbox{\dimexpr\linewidth-2\fboxsep\relax}{%
      \color{white}\bfseries\raggedright\strut#1\strut}}%
  \nopagebreak\par\nopagebreak
  \vspace{0.06cm}\nopagebreak
  \apresbandeautrue}
\newcommand{\partiefiche}[1]{%
  \refstepcounter{partiefiche}%
  \lycee@bandeau{\signet{1}{\thepartiefiche{} --- #1}\thepartiefiche{} --- #1}}
\newcommand{\partiefichesansnum}[1]{\lycee@bandeau{\signet{1}{#1}#1}}

\newcounter{exo}
\newlength{\espaceexo}\setlength{\espaceexo}{7pt}
\newlength{\espacetitre}\setlength{\espacetitre}{2pt}
\newcommand{\exo}[1][\relax]{%
  \par
  \ifapresbandeau\apresbandeaufalse\else\penalty-150\addvspace{\espaceexo}\fi
  \refstepcounter{exo}%
  \noindent\signet[exercice-\arabic{exo}]{2}{Exercice \arabic{exo}}%
  \textbf{\color{coulExo}\underline{Exercice \theexo}}%
  \ifx\relax#1\relax\else\textbf{ : #1}\fi
  \par\nopagebreak\vspace{\espacetitre}\nopagebreak
  \ignorespaces}

\setlist[enumerate]{nosep,leftmargin=*,itemsep=2pt,topsep=2pt}
\setlist[itemize]{nosep,leftmargin=*,topsep=2pt}
\newcommand{\choix}[3]{\par\hspace*{1em}%
  \makebox[0.3\linewidth][l]{a) #1}\makebox[0.3\linewidth][l]{b) #2}c) #3}
\newcommand{\deux}[2]{\par\hspace*{0.6em}%
  \makebox[0.48\linewidth][l]{#1}#2}
\newcommand{\trois}[3]{\par\hspace*{0.6em}%
  \makebox[0.32\linewidth][l]{#1}\makebox[0.32\linewidth][l]{#2}#3}

  \renewcommand{\exo}[1][\relax]{%
    \ifapresbandeau\apresbandeaufalse\else\par\penalty-150\fi
    \refstepcounter{exo}%
    \vspace{0.05cm}%
    \noindent\signet[exercice-\arabic{exo}]{2}{Exercice \arabic{exo}}%
    \textbf{\color{coulExo}\underline{Exercice \theexo}}%
    \ifx\relax#1\relax\else\textbf{ : #1}\fi%
    \par\nopagebreak
    \ignorespaces}
  \newcommand{\rappel}[1]{%
    \par\vspace{0.06cm}\nopagebreak
    \noindent\fcolorbox{black!35}{black!5}{%
      \parbox{\dimexpr\linewidth-2\fboxsep-2\fboxrule\relax}{%
        \small\textbf{Rappel.} #1}}%
    \par\vspace{0.06cm}\nopagebreak}
  \setlist[enumerate]{nosep,leftmargin=*}
  \setlist[itemize]{nosep,leftmargin=*}
  \setlength{\multicolsep}{3pt plus 1pt minus 1pt}
  \setlength{\premulticols}{10pt}
  \setlength{\postmulticols}{10pt}
  \newlist{expr}{enumerate}{1}
  \setlist[expr]{label={\Alph*\ $=$},nosep,leftmargin=*,itemsep=0pt}

\renewenvironment{center}{\par\addvspace{3pt}\centering}{\par\addvspace{3pt}}
\formulesserrees
\AtBeginDocument{\formulesserrees}
\clubpenalty=10000
\widowpenalty=10000
\displaywidowpenalty=10000
\brokenpenalty=10000
\setlength{\parskip}{0pt}

\RequirePackage{ragged2e}
\setlength{\RaggedRightRightskip}{0pt plus 1fil}
\AtBeginDocument{\RaggedRight\binoppenalty=10000 \relpenalty=10000 }
\g@addto@macro\@parboxrestore{\RaggedRight}
\makeatother
\usepackage{tkz-tab}

\newtcblisting{python}{%
  listing only,enhanced,
  colback=black!3,colframe=black!35,boxrule=0.5pt,arc=1pt,
  left=4pt,right=4pt,top=1pt,bottom=1pt,before skip=3pt,after skip=3pt,
  listing options={language=Python,basicstyle=\ttfamily\small,
    keywordstyle=\color{coulMeth}\bfseries,columns=fullflexible,
    keepspaces=true,showstringspaces=false,aboveskip=0pt,belowskip=0pt}}

\begin{document}

\titrefiche{1\textsuperscript{re} STI2D --- Primitives}{Fiche d'exercices du chapitre}

\begin{multicols}{2}\raggedcolumns

\partiefiche{Primitive d'une fonction}

\exo[Vérifier une primitive]
Dans chaque cas, dire si $F$ est une primitive de $f$.
\begin{enumerate}[label=\alph*)]
  \item $f(x)=5$ \quad et \quad $F(x)=5x-2$
  \item $f(x)=6x+1$ \quad et \quad $F(x)=3x^{2}+x-4$
  \item $f(x)=2x-7$ \quad et \quad $F(x)=x^{2}-7$
  \item $f(x)=3x^{2}-8x+2$ \quad et \quad $F(x)=x^{3}-4x^{2}+2x+9$
  \item $f(x)=4x-\dfrac{1}{x^{2}}$ \quad et \quad $F(x)=2x^{2}+\dfrac{1}{x}$ \quad ($x>0$)
  \item $f(x)=6x^{2}+\dfrac{5}{x^{2}}$ \quad et \quad $F(x)=2x^{3}+\dfrac{5}{x}$ \quad ($x>0$)
\end{enumerate}

\exo[En phrase complète]
Écrire chaque information par une phrase complète avec le mot « primitive » s'il
s'agit d'une dérivée, avec le mot « dérivée » s'il s'agit d'une primitive.
\begin{enumerate}[label=\arabic*)]
  \item La dérivée de $G(x)=x^{3}+4x$ est $g(x)=3x^{2}+4$.
  \item $H(x)=5x^{2}-x$ est une primitive de $h(x)=10x-1$.
  \item La vitesse $v$ d'un vélo est la dérivée de sa position $d$.
\end{enumerate}

\exo[Laquelle est une primitive ?]
Pour chaque fonction, trouver parmi les trois fonctions qui la suivent celle qui en est
une primitive, en justifiant.
\begin{enumerate}[label=\alph*)]
  \item $f(x)=4x-3$ : \quad $F_{1}(x)=4x^{2}-3x$ ; \quad $F_{2}(x)=2x^{2}-3x+1$ ;
        \quad $F_{3}(x)=4$
  \item $g(x)=3x^{2}+2$ : \quad $G_{1}(x)=6x$ ; \quad $G_{2}(x)=x^{3}+2$ ;
        \quad $G_{3}(x)=x^{3}+2x-5$
  \item $h(x)=-\dfrac{2}{x^{2}}$ \ ($x>0$) : \quad $H_{1}(x)=-\dfrac{2}{x}$ ;
        \quad $H_{2}(x)=\dfrac{2}{x}+3$ ; \quad $H_{3}(x)=-2x$
\end{enumerate}

\exo[Sur une piste d'essai]
Sur une piste d'essai, un vélo à assistance électrique démarre. Au bout de $t$ secondes
(pour $0\leq t\leq 8$), il a parcouru ${d(t)=0{,}25t^{2}+2t}$ mètres.
\begin{enumerate}[label=\arabic*)]
  \item Calculer la distance parcourue au bout de $4$ secondes.
  \item Montrer que $d$ est une primitive de la fonction $v$ définie par
        ${v(t)=0{,}5t+2}$.
  \item En physique, que représente $v(t)$ ?
  \item Calculer $v(8)$, puis le convertir en km/h.
\end{enumerate}

\partiefiche{Toutes les primitives d'une fonction}

\exo[Primitive qui vérifie une condition]
Pour chaque fonction $f$, définie sur $\R$, la fonction $F$ en est une primitive. Déterminer la primitive
$G$ de $f$ qui vérifie la condition écrite en dernier.
\begin{enumerate}[label=\alph*)]
  \item $f(x)=4x+1$ ; \quad $F(x)=2x^{2}+x+5$ ; \quad $G(1)=0$
  \item $f(x)=6x-2$ ; \quad $F(x)=3x^{2}-2x$ ; \quad $G(2)=1$
  \item $f(x)=3x^{2}-4x$ ; \quad $F(x)=x^{3}-2x^{2}+7$ ; \quad $G(-1)=2$
  \item $f(x)=5x^{4}+3x^{2}$ ; \quad $F(x)=x^{5}+x^{3}$ ; \quad $G(-1)=-3$
  \item $f(x)=x^{2}+1$ ; \quad $F(x)=\dfrac{x^{3}}{3}+x$ ; \quad $G(3)=4$
\end{enumerate}

\exo[Deux primitives]
Soit les fonctions $F$ et $G$ définies sur $\R$ par ${F(x)=x^{2}-6x+4}$ et ${G(x)=(x-3)^{2}}$.
\begin{enumerate}[label=\arabic*)]
  \item Développer $G(x)$.
  \item Calculer $F'(x)$ et $G'(x)$. Que peut-on en déduire pour $F$ et $G$ ?
  \item Calculer $G(x)-F(x)$. Quelle propriété du cours retrouve-t-on ?
  \item Déterminer la primitive $H$ sur $\R$ de ${f(x)=2x-6}$ telle que ${H(0)=0}$.
\end{enumerate}

\exo[Faire monter une courbe]
\begin{minipage}[c]{0.58\linewidth}\raggedright
Ci-contre, $\mathcal{C}_{F}$ est la courbe de la fonction $F$ définie par
${F(x)=x^{2}-2x-3}$.
\begin{enumerate}[label=\arabic*)]
  \item Vérifier que $F$ est une primitive de ${f(x)=2x-2}$.
  \item Lire $F(0)$ sur le graphique, puis le calculer.
  \item Déterminer la primitive $G$ de $f$ telle que ${G(0)=-1}$.
  \item Calculer $G(x)-F(x)$. En déduire comment obtenir la courbe de $G$ à partir de
        $\mathcal{C}_{F}$, puis la tracer.
\end{enumerate}
\end{minipage}\hfill
\begin{minipage}[c]{0.4\linewidth}\centering
\begin{tikzpicture}[x=0.4cm,y=0.4cm]
  \repere{-3}{-5}{5}{5}
  \gradx{-2,-1,1,2,3,4} \grady{-4,-3,-2,-1,1,2,3,4}
  \draw[coulCourbe,line width=1.1pt,domain=-1.6:3.6,samples=50]
    plot (\x,{\x*\x-2*\x-3});
  \node[coulCourbe,font=\footnotesize] at (4.2,2.4) {$\mathcal{C}_{F}$};
\end{tikzpicture}
\end{minipage}

\exo[Un ballon-sonde]
Un ballon-sonde météorologique est lâché du toit d'une station météo, à $15$~m du sol.
Il monte à la vitesse constante de $5$~m/s : au bout de $t$ secondes, sa vitesse est
${v(t)=5}$. Son altitude $h(t)$ (en m) est la primitive de $v$ telle que ${h(0)=15}$.
\begin{enumerate}[label=\arabic*)]
  \item Vérifier que ${F(t)=5t}$ est une primitive de $v$.
  \item Déterminer $h(t)$.
  \item Calculer l'altitude du ballon au bout d'une minute.
  \item Au bout de combien de temps le ballon atteint-il $1\,015$~m d'altitude ?
\end{enumerate}

\partiefiche{Primitive d'un polynôme}

\exo[Primitives de polynômes]
Déterminer une primitive de chaque fonction, définie sur $\R$.
\begin{multicols}{2}
\begin{enumerate}[label=\alph*)]
  \item $f(x)=7$
  \item $g(x)=6x$
  \item $h(x)=9x^{2}$
  \item $k(x)=4x^{3}-2x+3$
  \item $m(x)=12x^{2}-10x+1$
  \item $p(x)=5x^{4}+8x^{3}-6x^{2}$
  \item $q(x)=x^{2}+3x-4$
  \item $r(x)=-2x^{3}+5x-\dfrac{1}{3}$
\end{enumerate}
\end{multicols}

\exo[Toutes les primitives]
Déterminer toutes les primitives de chaque fonction, définie sur $\R$ (en c) et d), transformer d'abord
l'écriture).
\begin{multicols}{2}
\begin{enumerate}[label=\alph*)]
  \item $f(x)=2x+5$
  \item $g(x)=3x^{2}-4x+1$
  \item $h(x)=\dfrac{3x^{2}-2x}{2}$
  \item $k(x)=(x+1)(x-3)$
\end{enumerate}
\end{multicols}

\exo[Avec une condition initiale]
Pour chaque fonction, définie sur $\R$, déterminer sa primitive qui vérifie la condition écrite.
\begin{multicols}{2}
\begin{enumerate}[label=\alph*)]
  \item $f(x)=4$ ; \ $F(1)=0$
  \item $g(x)=2x-6$ ; \ $G(3)=1$
  \item $h(x)=6x^{2}+2x$ ; \ $H(1)=0$
  \item $k(x)=8x^{3}-3x^{2}+1$ ; \ $K(-1)=5$
  \item $m(x)=x^{2}-1$ ; \ $M(3)=0$
\end{enumerate}
\end{multicols}

\exo[Variations d'une primitive]
Soit la fonction $f$ définie sur $\R$ par $f(x)=3x^{2}-12$, et $F$ la primitive de $f$ telle que ${F(0)=5}$.
\begin{enumerate}[label=\arabic*)]
  \item Déterminer $F(x)$.
  \item Montrer que ${f(x)=3(x-2)(x+2)}$.
  \item Dresser le tableau de signes de $f(x)$ sur $\intff{-3}{3}$.
  \item Que vaut $F'(x)$ ? En déduire le tableau de variations de $F$ sur
        $\intff{-3}{3}$, avec les valeurs de $F$.
\end{enumerate}

\exo[Une fusée à eau]
Le club de sciences du lycée lance une fusée à eau à la verticale. À l'instant ${t=0}$,
toute l'eau est éjectée : la fusée est à $2$~m du sol et monte à $20$~m/s. Ensuite, la
pesanteur la freine : son accélération vaut ${a(t)=-10}$ (en m/s$^{2}$). La vitesse $v$
(en m/s) est une primitive de l'accélération, et l'altitude $h$ (en m) une primitive de
la vitesse.
\begin{enumerate}[label=\arabic*)]
  \item Déterminer $v(t)$, sachant que ${v(0)=20}$.
  \item Déterminer $h(t)$, sachant que ${h(0)=2}$.
  \item À quel instant la fusée cesse-t-elle de monter ?
  \item Calculer son altitude maximale.
  \item Montrer qu'au bout de $4$~s, la fusée est revenue à son altitude de départ.
\end{enumerate}

\partiefiche{Primitives de cosinus et sinus}

\exo[Primitives avec cosinus et sinus]
Déterminer une primitive de chaque fonction, définie sur $\R$.
\begin{multicols}{2}
\begin{enumerate}[label=\alph*)]
  \item $f(x)=5\cos(x)$
  \item $g(x)=-4\sin(x)$
  \item $h(x)=2\sin(x)+\cos(x)$
  \item $k(x)=3x^{2}-2\cos(x)$
  \item $m(x)=6x+7\sin(x)$
  \item $p(x)=1-\sin(x)$
\end{enumerate}
\end{multicols}

\exo[Avec une condition initiale]
Pour chaque fonction, définie sur $\R$, déterminer sa primitive qui vérifie la condition écrite.
\rappel{$\cos(0)=1$ ; \ $\sin(0)=0$ ; \ $\cos\left(\dfrac{\pi}{2}\right)=0$ ; \ $\sin\left(\dfrac{\pi}{2}\right)=1$ ;
\ $\cos(\pi)=-1$ ; \ $\sin(\pi)=0$.}
\begin{enumerate}[label=\alph*)]
  \item $f(x)=\cos(x)$ ; \ $F(0)=3$
  \item $g(x)=\sin(x)$ ; \ $G(0)=0$
  \item $h(x)=2\cos(x)+4x$ ; \ $H(0)=-1$
  \item $k(x)=3\sin(x)-\cos(x)$ ; \ $K(\pi)=2$
  \item $m(x)=-2\sin(x)$ ; \ $M\left(\dfrac{\pi}{2}\right)=4$
\end{enumerate}

\exo[Dériver pour trouver une primitive]
Soit la fonction $F$ définie sur $\R$ par $F(x)=x\sin(x)$.
\begin{enumerate}[label=\arabic*)]
  \item Calculer $F'(x)$ (dérivée d'un produit).
  \item En déduire une primitive de ${f(x)=\sin(x)+x\cos(x)}$.
  \item Déterminer la primitive $G$ de $f$ telle que ${G(\pi)=1}$.
  \item Soit la fonction $g$ définie sur $\R$ par $g(x)=x\cos(x)$. Écrire $g(x)$ à l'aide de $f(x)$ et de $\sin(x)$, puis
        en déduire une primitive de $g$.
\end{enumerate}

\exo[Vrai ou faux]
Dire si chaque affirmation est vraie ou fausse, en justifiant.
\begin{enumerate}[label=\arabic*)]
  \item $\sin(x)$ est une primitive de $\cos(x)$.
  \item $\cos(x)$ est une primitive de $\sin(x)$.
  \item ${F(x)=\sin(x)+2}$ et ${G(x)=\sin(x)-5}$ sont deux primitives d'une même fonction.
  \item La primitive de ${f(x)=-\sin(x)}$ qui vaut $0$ en $0$ est ${F(x)=\cos(x)}$.
  \item $\dfrac{x^{4}}{4}-7$ est une primitive de $x^{3}$.
\end{enumerate}

\partiefiche{Primitives de $A\cos(\omega t+\varphi)$ et $A\sin(\omega t+\varphi)$}

\exo[Primitives de fonctions sinusoïdales]
Déterminer une primitive de chaque fonction.
\begin{multicols}{2}
\begin{enumerate}[label=\alph*)]
  \item $f(t)=8\cos(4t)$
  \item $g(t)=9\sin(3t)$
  \item $h(t)=\cos(2t)$
  \item $k(t)=4\cos\left(2t+\dfrac{\pi}{3}\right)$
  \item $m(t)=10\sin\left(5t-\dfrac{\pi}{2}\right)$
  \item $p(t)=-6\cos(3t+\pi)$
\end{enumerate}
\end{multicols}

\exo[Toutes les primitives]
Déterminer toutes les primitives de chaque fonction.
\begin{multicols}{2}
\begin{enumerate}[label=\alph*)]
  \item $f(t)=4t-6\sin(2t)$
  \item $g(t)=5-8\cos\left(4t-\dfrac{\pi}{6}\right)$
  \item $h(t)=6t^{2}+2\cos(3t)$
  \item $k(t)=3\sin(2t)-4\sin(4t)$
\end{enumerate}
\end{multicols}

\exo[Avec une condition initiale]
Pour chaque fonction, déterminer sa primitive qui vérifie la condition écrite.
\begin{enumerate}[label=\alph*)]
  \item $f(t)=6\cos(3t)$ ; \ $F(0)=1$
  \item $g(t)=4\sin(2t)$ ; \ $G(0)=0$
  \item $h(t)=4\cos(2t)$ ; \ $H\left(\dfrac{\pi}{4}\right)=3$
  \item $k(t)=6\sin\left(3t+\dfrac{\pi}{2}\right)$ ; \ $K(0)=5$
\end{enumerate}

\exo[Courant et charge d'un condensateur]
Dans un circuit, un condensateur reçoit un courant d'intensité ${i(t)=2\cos(100t)}$ (en
ampères, $t$ en secondes). Sa charge $q(t)$ (en coulombs) vérifie ${q'(t)=i(t)}$ et
${q(0)=0}$.
\begin{enumerate}[label=\arabic*)]
  \item Compléter la phrase avec un mot du cours : « $q$ est une \dots{} de $i$ ».
  \item Déterminer $q(t)$.
  \item Calculer $q\left(\dfrac{\pi}{200}\right)$.
  \item Vérifier le résultat de la question 2 en dérivant $q$.
\end{enumerate}

\exo[L'énergie des vagues]
Un houlomoteur produit de l'électricité avec les vagues : son flotteur monte et descend
avec la houle. Au bout de $t$ secondes, la vitesse verticale du flotteur (en m/s) est
${v(t)=\sin(0{,}5t)}$, et sa hauteur au-dessus du fond (en m) est la primitive $h$ de
$v$ telle que ${h(0)=8}$.
\begin{enumerate}[label=\arabic*)]
  \item Déterminer $h(t)$.
  \item Calculer $h(2\pi)$.
  \item Calculer la période de $h$ (au chapitre 7, ${T=\dfrac{2\pi}{\omega}}$),
        arrondie au dixième de seconde.
  \item Justifier que le flotteur reste entre $8$~m et $12$~m de hauteur.
\end{enumerate}

\partiefiche{Méthode d'Euler}

\exo[Méthode d'Euler]
Soit ${f(t)=\dfrac{6}{t+2}}$ pour $t\geq 0$, et $F$ la primitive de $f$ telle que
${F(0)=0}$.
\begin{enumerate}[label=\arabic*)]
  \item Avec la méthode d'Euler et un pas ${h=0{,}5}$, compléter le tableau (valeurs
        arrondies à $0{,}01$).
\end{enumerate}
\begin{center}
\renewcommand{\arraystretch}{1.4}\setlength{\tabcolsep}{5pt}
\begin{tabular}{|c|c|c|c|c|c|}
\hline
$t$ & $0$ & $0{,}5$ & $1$ & $1{,}5$ & $2$\\
\hline
$f(t)$ & & & & & \\
\hline
$F(t)\approx$ & $0$ & & & & \\
\hline
\end{tabular}
\end{center}
\begin{enumerate}[label=\arabic*),start=2]
  \item Placer les points dans un repère (unité $2$~cm), puis les relier.
\end{enumerate}

\exo[Un pas plus petit]
Soit ${f(t)=\dfrac{2}{t+1}}$ pour $t\geq 0$, et $F$ la primitive de $f$ telle que
${F(0)=0}$.
\begin{enumerate}[label=\arabic*)]
  \item Avec la méthode d'Euler et un pas ${h=0{,}5}$, calculer une valeur approchée de
        $F(1)$, arrondie à $0{,}01$.
  \item Même question avec un pas ${h=0{,}25}$ (quatre étapes).
  \item Avec un pas très petit, un programme donne ${F(1)\approx 1{,}39}$. Quel pas
        donne la meilleure approximation ?
\end{enumerate}

\exo[Programmer la méthode d'Euler]
Le programme ci-dessous, incomplet, applique la méthode d'Euler à
${f(t)=\dfrac{4}{1+t^{2}}}$, avec ${F(0)=0}$.
\begin{python}
def f(t):
    return .....

def euler(t, y, h, n):
    for i in range(n):
        y = .....
        t = .....
    return y
\end{python}
\begin{enumerate}[label=\arabic*)]
  \item Compléter les trois lignes en pointillés.
  \item Pour l'appel \texttt{euler(0, 0, 0.5, 2)}, donner à la main les valeurs de
        \texttt{t} et de \texttt{y} après chaque tour de boucle.
  \item Quel appel calcule une valeur approchée de $F(1)$ avec un pas de $0{,}01$ ?
\end{enumerate}

\exo[Un bateau électrique ralentit]
Un bateau électrique coupe son moteur et glisse sur un lac en ralentissant. Au bout de
$t$ secondes, sa vitesse (en m/s) est ${v(t)=\dfrac{12}{t+2}}$. La distance parcourue
$d(t)$ (en m) vérifie ${d'(t)=v(t)}$ et ${d(0)=0}$.
\begin{enumerate}[label=\arabic*)]
  \item Calculer la vitesse du bateau au départ, puis au bout de $4$~s.
  \item Montrer que ${v'(t)=-\dfrac{12}{(t+2)^{2}}}$, puis justifier que le bateau
        ralentit.
  \item Avec la méthode d'Euler et un pas ${h=1}$, calculer des valeurs approchées de
        $d(1)$ ; $d(2)$ ; $d(3)$ et $d(4)$.
  \item Placer les points dans un repère, puis les relier.
\end{enumerate}

\end{multicols}

\end{document}
